% s-120-p120-m002; original PDF 120, 121
\Aphorism{s-120-p120-m002}{150}
Wenn man ein Weib liebt, ist man meilenweit davon, es zu kennen. Man verlässt sich auf Inspiration, statt den analytischen Geduldsweg zu den Einzelheiten zu gehen; man stellt das Ideal des Weibes auf, im Ganzen und Ungefähren, als nothwendige Hypothese, als Gott, als einzige Rettung aus einzigen Nöthen, man hat nichts als ein vages schwimmendes tumultuarisches Gemeingefühl an Stelle der zahllosen Einzelreactionen, die sonst das Bild einer Person zusammensetzen. Der object-schaffende Trieb, die kühle zeichnerische Sachlichkeit des Bewusstseins versagen an der Geliebten; es ist kein Raum in der Seele, ihr Bild aufzurichten, keine Fläche, sie darauf abzuschatten, keine Ferne, sie zu betrachten.
\par

Ein Weib, das uns als Grundempfindung ausfüllt, verbannt sich ebendamit als Gestalt und Porträt aus unserer Phantasie; späterhin, nach erloschener Liebe, sind wir erstaunt, zum ersten Male dahinter zu kommen, wie die Geliebte eigentlich aussah.
\par

